\documentclass[10pt,a4paper]{amsart}
\usepackage[margin=19mm]{geometry}
\usepackage{amsmath,amssymb,mathtools}
\usepackage[scr=rsfs,cal=euler]{mathalfa}
\usepackage{xfrac}
\usepackage[unicode,hidelinks,bookmarks=false]{hyperref}
\newcommand{\ie}{i.e.\ }
\newcommand{\cf}{cf.\ }
\newcommand{\resp}{respectively\ }
\newcommand{\Subsection}{\subsection}
\providecommand{\nobreakdash}{\nobreak\hbox{-}}
\setlength{\emergencystretch}{2em}
\multlinegap0pt
\usepackage{mathtools}
\usepackage{amssymb}
\usepackage[shortlabels]{enumitem}
\newtheorem{lem}{Lemma}
\usepackage{cleveref}
\crefname{lem}{Lemma}{Lemmas}
\title{subgroups with all finite preimages isomorphic are conjugate}
\author{Ido Karshon, Alexander Lubotzky, D. B. McReynolds, Alan W. Reid, Mark Shusterman}
\date{Source: arXiv:2602.15463v3 (CC BY 4.0)}
\begin{document}
\maketitle

\noindent\emph{Source and licence.} subgroups with all finite preimages isomorphic are conjugate. Version arXiv:2602.15463v3 (CC BY 4.0), distributed by arXiv under the Creative Commons Attribution 4.0 International licence, http://creativecommons.org/licenses/by/4.0/, as recorded in the arXiv licence metadata for this exact version and shown on the abstract page https://arxiv.org/abs/2602.15463v3. Full licence notice: \texttt{licenses/arxiv-2602.15463-CC-BY-4.0.txt}.\par\smallskip

\noindent\textit{Editorial scope.} Counted unit: the article's lem lem:rigidity (source0.tex lines 316--318). The article's complete original proof of it is delivered with it (source0.tex lines 320--329, one \texttt{proof} environment of 10 lines). No mathematics is written, repaired or re-derived: the statement and the proof are the article's own lines, in the article's own notation, and no retained line is substituted or re-flowed.  Retained uncounted as the source-local context the proof needs: lines 300--302; lines 308--310.  Nothing was omitted from any retained span, so the package carries no omission record.  No retained line contains a citation, so the wrapper carries no bibliography and no bibliography entry.  No retained line contains a figure environment, an image-inclusion command or drawing code, so no image is delivered and the package declares no asset dependency.  Editorial formatting only, in addition: an AMS \texttt{amsart} preamble that restates the article's own declarations and macros used by the retained lines, namely the environments \texttt{lem} on separate counters, together with the standard packages those lines need.  The retained environments print in this document as Lemma 1 in order of appearance, whereas the article prints the same items with the numbering of its own full document.  The complete native source of the article as distributed in the arXiv e-print for this exact version is bundled unchanged as \texttt{sources/194799/source0.tex}.
\begin{lem} \label{TrivCentInn}
The subgroup $\mathrm{Inn}(S)$ of $\mathrm{Aut}(S)$ is normal and its centralizer is trivial.
\end{lem}

\begin{lem}\label{lem:centralizer-and-minimal-normal}
Let $P$ be a subgroup of $\mathrm{Aut}(S)$ that contains $\mathrm{Inn}(S)$. Then every nontrivial normal subgroup $N$ of $P$ contains $\mathrm{Inn}(S)$. Consequently $\mathrm{Inn}(S)$ is the unique minimal nontrivial normal subgroup of $P$.
\end{lem}

\begin{lem} \label{lem:rigidity}
Let $P,Q$ be subgroups of $\mathrm{Aut}(S)$ containing $\mathrm{Inn}(S)$, and let $\Psi \colon P\xrightarrow{}Q$ be a group isomorphism. Then there exists $a \in \mathrm{Aut}(S)$ such that $\Psi(p)=a p a^{-1}$ for every $p\in P$.
\end{lem}

\begin{proof}
It follows from \cref{lem:centralizer-and-minimal-normal} applied to $P$ and to $Q$ that $\Psi(I)=I$, so
\[ \Psi|_{\mathrm{Inn}(S)} \colon \mathrm{Inn}(S) \xrightarrow{} \mathrm{Inn}(S) \]
is an automorphism, hence it gives rise to an automorphism of $S$, which we denote by $a \in \mathrm{Aut}(S)$. That is, for every $s \in S$ we have $\Psi(c_s)=c_{a(s)}=a c_s a^{-1}$ so for every $i \in \mathrm{Inn}(S)$ we get that
\[ \Psi(i)=aia^{-1}. \]

We want to show that $\Psi(p) = apa^{-1}$, or equivalently $\Psi(p)^{-1} a p a^{-1} = \mathrm{id}_S$, for each $p \in P$. By \cref{TrivCentInn} it suffices to show that $\Psi(p)^{-1} a p a^{-1}$ commutes with every $i \in \mathrm{Inn}(S)$, or equivalently
\[ (a p a^{-1}) i (a p a^{-1})^{-1} = \Psi(p) i \Psi(p)^{-1}. \]
Indeed we have $a p a^{-1} i a p^{-1}a^{-1} = \Psi(p a^{-1} i a p^{-1}) = \Psi(p)\Psi(a^{-1} i a)\Psi(p)^{-1} = \Psi(p) i \Psi(p)^{-1}$.
\end{proof}

\end{document}
